\documentclass[11pt]{article}
\usepackage[letterpaper,margin=0.82in]{geometry}
\usepackage{amsmath,amssymb,amsthm,mathtools,booktabs,longtable,array}
\usepackage[T1]{fontenc}
\usepackage{lmodern,microtype,xcolor,xurl,hyperref,fancyhdr}
\hypersetup{colorlinks=true,linkcolor=blue!45!black,urlcolor=blue!45!black,citecolor=blue!45!black,
 pdftitle={Synchronized exponent certificates for ZFC consistency},
 pdfauthor={AI-assisted research prepared for Vladimir Reshetnikov},
 pdfsubject={An explicit 358-variable natural quartic and exact evaluator reductions}}
\pagestyle{fancy}\fancyhf{}\fancyhead[L]{\small Diophantine consistency}
\fancyhead[R]{\small 24 October 2026}\fancyfoot[C]{\thepage}
\setlength{\headheight}{14pt}\setlength{\parskip}{4pt}
\newtheorem{theorem}{Theorem}[section]
\newtheorem{lemma}[theorem]{Lemma}
\newtheorem{proposition}[theorem]{Proposition}
\newtheorem{corollary}[theorem]{Corollary}
\theoremstyle{definition}\newtheorem{definition}[theorem]{Definition}
\newcommand{\N}{\mathbb N}\newcommand{\Z}{\mathbb Z}\newcommand{\Q}{\mathbb Q}
\newcommand{\Con}{\operatorname{Con}}\newcommand{\Refut}{\operatorname{Ref}}
\newcommand{\file}[1]{\nolinkurl{#1}}
\newcommand{\ceil}[1]{\left\lceil#1\right\rceil}
\newcommand{\floor}[1]{\left\lfloor#1\right\rfloor}
\newcommand{\Pmain}{P_{358}}
\allowdisplaybreaks[2]
\title{Synchronized exponent certificates\\for ZFC consistency\\[5pt]
\large An explicit quartic in 358 nonnegative variables,\\with exact contractions of five smaller-variable evaluators}
\author{Research prepared for Vladimir Reshetnikov\\
\small AI-assisted mathematical construction and separate arithmetic reviews}
\date{24 October 2026}

\begin{document}
\maketitle
\begin{abstract}
We give a completely specified polynomial $\Pmain\in\Z[X_0,\ldots,X_{357}]$
of exact degree four whose lack of nonnegative integer zeros is provably
equivalent in ZFC to $\Con(\mathrm{ZFC})$. The polynomial is the sum of
the squares of 275 explicitly printed quadratic residuals. Its collected
expansion has 8,229 monomials and coefficient height
$9\,903\,520\,314\,287\,616\,167\,992\,362\,844$.
The construction uses the repository's direct numerical C222 proof
calculus, without a K215 execution encoding or Jones universal polynomial.

The main change is to give two large arithmetic quantities the same
exponent so that they share an indexed Pell certificate. A common
factorial radius and an additional Pell index, introduced solely to serve
two addition identities, reduce the source further. The best direct
natural quartic found in the inspected incoming material had 395 variables;
the new equation removes 37, or about 9.37 percent. Its larger coefficient
height is a real cost. A 361-variable companion has fewer monomials.

Separately, exact ring identities reduce five positive-domain evaluators
to 98, 103, 107, 95 and 87 binary operations at their unchanged degrees
368, 340, 211, 70 and 51. Complete coefficient data, generators, independent
checks and proof dependencies accompany this report. The proofs are
ordinary mathematical arguments formalizable in ZFC; no new Lean proof
or encoded object-language ZFC derivation of the final polynomial is claimed.
\end{abstract}
\clearpage
\tableofcontents
\newpage

\section{The result and what is new}
Throughout, $\N=\{0,1,2,\ldots\}$. A positive domain is written explicitly
as $\Z_{>0}$; it is not interchangeable with the unrestricted integer domain.
We use the ordinary syntactic consistency statement
\[
 \Con(\mathrm{ZFC})\equiv
 \neg\exists p\in\N\;\operatorname{Proof}_{\mathrm{ZFC}}(p,\bot).
\]
The construction does not assume this statement or its negation.

\begin{theorem}[Explicit consistency equation]\label{thm:main}
Let $r_0,\ldots,r_{274}$ be exactly the integer polynomials in
Appendix~\ref{app:residuals}, and put
\begin{equation}\label{eq:main}
 \boxed{\displaystyle \Pmain(X_0,\ldots,X_{357})
       =\sum_{j=0}^{274}r_j(X_0,\ldots,X_{357})^2.}
\end{equation}
Then the ordinary arguments in this report, composed with the pinned
C222/refutation calibration, establish
\begin{equation}\label{eq:zfc}
 \mathrm{ZFC}\vdash
 \left[\neg\exists\boldsymbol X\in\N^{358}\;
       \Pmain(\boldsymbol X)=0\right]
 \longleftrightarrow\Con(\mathrm{ZFC}).
\end{equation}
The polynomial has exact degree four and 8,229 nonzero collected monomials.
Its coefficient height is
\[
 9\,903\,520\,314\,287\,616\,167\,992\,362\,844.
\]
\end{theorem}

The full residual table is printed rather than left as a request to run a
compiler. The identical table is stored in
\file{direct/natural358/equation.json}; the complete collected expansion is
\file{direct/natural358/polynomial.json}. Indices in these files are zero
based, just as in the appendix. Every coefficient is a literal integer.
No variable exponent, factorial, binomial coefficient, bounded quantifier,
universal-polynomial index or unspecified numerical parameter remains in
\eqref{eq:main}. Later uses of such operations explain how its finite
quadratic certificate was designed and why it is correct.

\subsection{Comparison at a fixed repository snapshot}
The repository snapshot is
\file{84e16e43092614bae262ca59651cb9d77e31a255}. The comparison includes
the active research and the incoming archives, not only the numbered
manuscript collection. The active comparison is materially ahead of that
collection's introductory summary~\cite{repo,frontier}.

\begin{center}\small
\begin{tabular}{lrrrr}
\toprule
Direct natural construction & Variables & Residuals & Terms & Height bits\\
\midrule
Active candidate-output radii & 399 & 317 & 8,286 & 27\\
Incoming centered factorials & 395 & 313 & 8,292 & 27\\
Shared-square integration alone & 395 & 313 & 8,131 & 27\\
New, smaller support & 361 & 278 & 6,759 & 94\\
\textbf{New, fewer variables} & \textbf{358} & \textbf{275} & \textbf{8,229} & \textbf{94}\\
\bottomrule
\end{tabular}
\end{center}
Every equation in this table is quartic over $\N$. The 395-variable
shared-square point is a measured intermediate, not another new variable
record. The two final variants are included in the deliverable. Their
coefficient heights are given exactly in \file{direct/CENSUS.json}.

The 358-variable result improves the direct 395-variable comparison by
37 variables and the active 399-variable comparison by 41. It does not
improve every metric: its largest coefficient requires 94 bits, whereas
the smaller-height predecessors require 27. The existing 604-variable
\emph{unrestricted integer} quartic is a different-domain comparison.
The repository's nine-positive-variable equations and ten-integer
equations remain much smaller in variable count, with other degree and
coefficient tradeoffs. No global minimum of any unrestricted class is proved.

\subsection{Attribution and novelty boundary}
The incoming Bootstrap package already contains the squared-output
factorial fixed-point theorem and the idea of sharing $(H!)^2$ as a
radius~\cite{bootstrap}. The incoming Valuations and Centered Certificates
package already contains the centered factorial identity and its
395-variable quartic~\cite{centered}. These ideas were independently
rediscovered during this investigation and then recognized in the
incoming material; they are not claimed as new here.

The concrete contribution of this report is the composition of:
\begin{enumerate}
\item a carrier and a majorant chosen to share one exponent;
\item the existing common square radius with a uniform low-digit bound;
\item a genuinely certified extra Pell index, used only to derive two
      requested indices by addition;
\item exact natural-domain projections giving the delivered 358-variable
      polynomial.
\end{enumerate}
This combination and literal equation were not found in the inspected
snapshot. That is a repository-relative novelty claim, not a worldwide
priority claim. The five evaluator contractions in the supplement are
separate exact algebraic improvements.

\section{The direct proof-certificate source}
The source is the fixed numerical calculus C222. Its truth label is zero,
its falsity label is one, and its formula constructors start with
\[
 I(a,b)=2(a+b)(a+b+1)+4b.
\]
The membership, universal-quantifier and equality constructors use offsets
one, two and three. The source has twelve padded single-conclusion rule
branches: truth, eight unconditional schema families, the guarded bundle,
modus ponens and generalization. Unused premise slots are padded by zero.
The exact constructors and rule data are included under
\file{upstream/engine/upstream/provenance/upstream/}.

The logical calibration is supplied by the repository's actual
unconditional interfaces, rather than an axiomhood oracle:
\begin{align*}
 \mathrm{C222}(1)&\longleftrightarrow
       \operatorname{BProv}(\mathrm{ZF},[],\bot),\\
 \operatorname{BProv}(\mathrm{ZF},[],\bot)
       &\longleftrightarrow\operatorname{BProv}(\mathrm{ZFC},[],\bot).
\end{align*}
Their declarations are
\file{KEndpoint.c222_iff_zf_refutation} and
\file{ConstructibleChoice.zf_refutation_iff_zfc_refutation}.
The supplied calibration pins identify the exact Lean source copies.
The ordinary logical transport and its scope were reviewed in the pinned
direct-C222 report~\cite{directreview}. No machine execution is part of
the new polynomial's source.

\subsection{Why a multiset of rules suffices}
For a valid two-premise rule $(p,q)\to v$, define its flow vector as
$e_v-e_p-e_q$, with $e_0=0$. A proof tree of falsity gives a finite
multiset whose total flow is $e_1$. Conversely, suppose such a multiset
exists. Let $D$ be the least rule-closed set of derivable labels. If one
were outside $D$, sum all flow coordinates outside $D$. Any rule whose
conclusion is outside $D$ has at least one premise outside $D$, so its
contribution to that sum is at most zero. A rule with conclusion inside
$D$ also contributes at most zero. The target flow $e_1$ contributes one,
a contradiction. Only finitely many supported coordinates are summed.
Cycles and repeated premises therefore create no false refutation.

Choose a horizon $t$ at least two, at least the proof-tree length, and at
least every local schema parameter. Pad with truth rows, which have zero
flow. This justifies the new convention $t=t_0+2$, $t_0\in\N$ without
losing any finite refutation.

\subsection{One numerical carrier encodes all label balances}
The fixed source supplies a 266-gate label-majorant certificate. Its
computed exponent is
\[
 d_L=3\,445\,908,\qquad \lambda=2^{22}=4\,194\,304\ge d_L.
\]
For parameters at most $t$, every relevant label plus one is at most
$(t+1)^\lambda$. This is proved gate by gate: an input has weight one,
a constant $c>1$ has weight $\lceil\log_2c\rceil$, addition sends weights
$(a,b)$ to $\max(a,b)+1$, and multiplication sends them to $a+b$.

Put $M=(t+1)^{\lambda t}=M_0+1$. For a proposed list of rule instances,
let
\[
 A(Z)=\prod_{i<t}(1+v_iZ),\qquad
 D(Z)=(1+Z)\prod_{i<t}(1+p_iZ)(1+q_iZ).
\]
Choose $h\ge1$ bounding all labels, including the target label one.
Then $M\ge(h+1)^t$.
\begin{lemma}[One-sided degree bootstrap]\label{lem:carrier}
If $A(M)=D(M)$ with $M>1$ as above, then $A=D$ as integer polynomials,
and the rule instances have total flow $e_1$.
\end{lemma}
\begin{proof}
The upper bound
\[
 A(M)\le(1+hM)^t<((h+1)M)^t\le M^{t+1}
\]
shows that $D$ cannot have degree greater than $t$: a nonzero leading
coefficient is a positive integer. Thus both products have at most $t$
nonconstant factors. Their values at one are at most $(h+1)^t\le M$;
since their constant coefficients are one, every coefficient is strictly
smaller than $M$. Uniqueness of base-$M$ expansion implies $A=D$.
Multiplicities of the distinct roots $-1/j$, for $j>0$, give precisely
the positive-label flow balances.
\end{proof}

The source uses one accumulator,
\[
 S_0=M+1,\quad S_t=1,\quad
 S_{i+1}(1+Mv_i)=S_i(1+Mp_i)(1+Mq_i).
\]
All states are positive by induction from the initial endpoint, so
telescoping gives the carrier equality. Conversely, a proof-tree traversal
lets $S_i$ be the product of its current open-obligation factors. This
provides an integral accumulator rather than merely a rational path.

The accumulator is beta-coded by $x,y$ using the moduli $1+(i+1)y$.
The eight bounded local columns are the four schema parameters and
$(S,S',u,v_q)$, with respective bounds
\[
 (t,t,t,t,y,y,x,x).
\]
The endpoint equations and the complete local source are retained from
the accepted projected construction~\cite{parent407,parent399}.

\section{The main change: synchronize two exponents}
The earlier construction separately computed
\[
 M=((t+1)^t)^{2^{22}},\qquad H=R^{2^{23}},\qquad
 R=2+t+x+y+M_0.
\]
The outer powers were represented by 22 and 23 repeated squaring steps.
The construction already had a common Pell representation of arbitrary
positive powers. Merely moving both fixed powers to that representation
removes many wires, but initially creates two separate exponent tests.

The better choice is
\begin{equation}\label{eq:sharedexp}
 \boxed{t=t_0+2,\quad e=2^{22}t,\quad
 M=(t+1)^e=M_0+1,\quad H=R^e.}
\end{equation}
The carrier is unchanged at a given $t$. The majorant may be larger,
and that is exactly the freedom we use.

The local-source ledger has 338 gates and computed exponent
\[
 d_H=6\,893\,253<2^{23}.
\]
It bounds the source on its finite boxes, plus the local bounds and a
strict positive margin, by $R^{2^{23}}$. Since $t\ge2$,
\[
 e\ge2^{23},\qquad H\ge R^{2^{23}}\ge R^4.
\]
All subsequent uses of $H$ require lower bounds of this kind. None
requires an upper bound on $H$ or its old exact value. The factorial
$H!$ and every later modulus are computed for the new value of $H$.

Both requests in \eqref{eq:sharedexp} now use the same exponent atom and
therefore the same indexed Pell pair. This is an existence-preserving
change of the complete arithmetic presentation, not an identity between
the old and new final polynomials.

\begin{proposition}[General synchronization principle]
Suppose a carrier needs exponent $Lt$, a local bound needs $H\ge R^K$,
and the proof horizon can be padded to $t\ge q$. If $qL\ge K$, choosing
both carrier and majorant with exponent $Lt$ preserves the required
bounds and allows one indexed Pell certificate to serve both requests.
\end{proposition}
\begin{proof}
The carrier is unchanged, and $Lt\ge qL\ge K$ gives $R^{Lt}\ge R^K$
for $R\ge1$. Completeness chooses a padded horizon before either value.
The shared certificate is complete because its Pell pair is uniquely
determined by the common Pell parameter $A$ and index $e$.
\end{proof}
A smaller shared coefficient would have to satisfy
$L\ge\max(d_L,\lceil d_H/2\rceil)=3\,446\,627$ at the same horizon
offset. Simply taking $L=d_L$ would fail this check. The delivered
equation retains $2^{22}$, which keeps the source comparison transparent.

\section{Shared power certificates and the necessary proof order}
For $A\ge2$, define the Pell pairs by
\[
 X_0(A)=1,\quad Y_0(A)=0,\quad
 X_{k+1}=AX_k+(A^2-1)Y_k,\quad Y_{k+1}=X_k+AY_k.
\]
All recurrences here explain fixed finite polynomial certificates; they
are not extra clauses in the delivered equation.

\subsection{A frame independent of intended output meanings}
For the finite collection of distinct atoms occurring as bases, exponents
or proposed outputs of power requests, impose
\begin{equation}\label{eq:frame}
 W=2+\sum\text{atoms},\qquad
 A^2-W^3(W+2)Z^2=1,\qquad Z=z+1.
\end{equation}
Every coordinate is nonnegative. Thus $W\ge2$ and all those atoms are
strictly less than $W$. Rewrite the norm as
\[
 A^2-((W+1)^2-1)(WZ)^2=1.
\]
Pell classification gives $A=X_j(W+1)$ and $WZ=Y_j(W+1)$ for $j>0$.
The recurrence modulo $W$ gives $Y_j(W+1)\equiv j\pmod W$, hence
$W\mid j$ and $j\ge W$. Consequently
\begin{equation}\label{eq:framebound}
 A\ge(W+1)^j\ge(W+1)^W>W^W.
\end{equation}
This argument does not use the meaning of $H$, a factorial, a binomial
coefficient or any proposed power output.

For a positive base $n$, the literal modulus equation is
\[
 2An=T_n+n^2+1.
\]
It independently excludes $n=0$ over $\N$. For $1\le n<W$ it gives
$T_n\ge2A-2$. If a request has a positive exponent $k<W$ and proposed
output $m<W$, then both $m$ and $n^k$ are strictly smaller than $T_n$.

\subsection{One certificate for each seed index}
For a seed exponent $k>0$, the included Matiyasevich index certificate
identifies the nonnegative pair $(X,Y)$ with $(X_k(A),Y_k(A))$.
In its shared form it uses a common auxiliary Pell pair $(U,V)$ and base
$B_*$. Its three norms are
\begin{align*}
 X^2-(A^2-1)Y^2&=1,\\
 U^2-(A^2-1)V^2&=1,\\
 S^2-(B_*^2-1)R_*^2&=1.
\end{align*}
The actual common-base equations are
\[
 B_*=1+4V(q_4+1)=A+Uq_B,\qquad q_4,q_B\in\N.
\]
The certificate requires $Y\ge k$ and
\[
 B_*\equiv A\pmod U,\quad B_*\equiv1\pmod{4Y},\quad
 S\equiv X\pmod U,\quad R_*\equiv k\pmod{4Y},
\]
and $Y^2\mid V$. All quotients have explicit natural-coordinate forms.
The supplied inherited proof establishes simultaneous existence by finite
unit periods and compatible congruences. The implementation enforces the
stronger common divisibility
\begin{equation}\label{eq:seedproduct}
 V=\left(\prod_{k\text{ a seed}}Y_k^2\right)(v+1).
\end{equation}
Every actual seed is included in this product.
If $V=0$, its norm gives $U=1$, while the common-base equations give
$1=A+q_B$, contradicting $A\ge2$. Hence $V>0$, and
\eqref{eq:seedproduct} forces every seed ordinate to be positive,
including those whose exponent meanings have not yet been recovered.
Thus the selected known-positive exponent can be interpreted first.

The power row is
\begin{equation}\label{eq:power}
 X_k(A)+nY_k(A)=AY_k(A)+m+T_n h,\qquad h\in\N.
\end{equation}
Set $L_k=X_k(A)+(n-A)Y_k(A)$. The exact recurrence
\[
 L_0=1,\qquad L_{k+1}=nL_k+T_nY_k(A)
\]
gives $L_k\equiv n^k\pmod{T_n}$. The strict bounds therefore force
$m=n^k$. Conversely
\[
 L_k-n^k=T_n\sum_{j=0}^{k-1}n^{k-1-j}Y_j(A)
\]
supplies a nonnegative quotient. These are pointwise statements for an
individual request with known positive exponent.

\subsection{Why the new output \texorpdfstring{$H$}{H} is handled first}
Previously $H>0$ followed from its literal squaring syntax. It is now a
raw proposed power output, so assuming its positivity at the beginning
would be a gap. The valid order is:
\begin{enumerate}
\item Establish \eqref{eq:framebound} from the actual nonnegative atoms.
\item Apply the individual-request argument to the two calls in
      \eqref{eq:sharedexp}. Their bases $t+1,R$ and their common exponent
      $e$ are positive independently of all proposed outputs.
\item Recover $M=(t+1)^e$ and $H=R^e>0$. Now the old local bound is valid.
\item Exclude a zero proposed factorial $T$ by the base-$T$ modulus row
      of the actual request $T^t$; this uses no factorial interpretation.
\item Recover all remaining positive exponents, including the extra
      index introduced in the next section, and interpret all powers.
\end{enumerate}
Sharing the auxiliary frame does not force these semantic arguments to be
applied simultaneously. Its equations can be used first for a selected
known-positive index, then for the remaining indices.

\section{One square radius and an index used only for sharing}
Let $T$ and $J$ initially be proposed values of $H!$ and $t!$. Put
\begin{equation}\label{eq:rho}
 \rho=T^2,\qquad
 V_n=(\rho+1)^n,\quad C_n=\binom{\rho+n}{n}\quad(n=H,t).
\end{equation}
For $F_H=T$ and $F_t=J$, impose the two ceiling rows
\begin{equation}\label{eq:ceiling}
 V_n+C_n=C_nF_n+r_n+1,\qquad C_n=r_n+s_n+1,
 \qquad r_n,s_n\in\N.
\end{equation}
When $V_n,C_n>0$, these are exactly
$F_n=\lceil V_n/C_n\rceil$. In particular, they exclude $F_n=0$.

\subsection{The extra index at \texorpdfstring{$\rho$}{rho}}
The binomial graphs request powers with exponents $H+\rho$ and
$t+\rho$. Pairs at $H$ and $t$ already exist. We introduce a certified
pair at $\rho$ even though no arithmetic power request uses $\rho$
alone. Then use
\begin{align}
 X_{a+b}&=X_aX_b+(A^2-1)Y_aY_b,\\
 Y_{a+b}&=X_aY_b+Y_aX_b.\label{eq:addition}
\end{align}
One extra seed replaces two composite seeds, with two smaller addition
circuits. The third derived index is the already available $t+1$.

This additional seed is not free or unchecked. Its norm, index test and
squared ordinate in \eqref{eq:seedproduct} are all present in the
literal source. Its positivity follows before factorial recovery:
the base-$T$ modulus row gives $T\ge1$, hence $\rho\ge1$. Moreover
$\rho+1$ is an actual power base, so $\rho<W$ follows from the frame.
The finite-unit existence proof applies to the enlarged finite seed set.
The three addition dependencies are acyclic.

The final inventory contains 22 power requests, seven distinct bases,
18 indexed pairs, 15 certified seeds and three derived pairs. Every
identity in \eqref{eq:addition}, its exponent binding, and the complete
15-factor product are independently checked against the actual raw rows.

\subsection{A coefficient bound that avoids a factorial assumption}
\begin{lemma}[Uniform low-digit envelope]\label{lem:binomial}
For integers $\rho,n\ge1$ and $0\le j\le n$,
\[
 \binom{\rho+n}{j}\le(\rho+1)^n.
\]
\end{lemma}
\begin{proof}
Induct on $n$. At $n=1$ the values are one and $\rho+1$.
For $0<j<n$, Pascal's identity and the induction hypothesis give
$2(\rho+1)^{n-1}\le(\rho+1)^n$; $j=0$ is immediate. For $j=n$ use
\[
 \binom{\rho+n}{n}=\prod_{k=1}^n\frac{\rho+k}{k}
 \le\prod_{k=1}^n(\rho+1).
\]
\end{proof}
No assumption $\rho\ge n$ is needed. This is essential while $T$ is
still only a proposed factorial. The hypothesis $\rho\ge1$ cannot be
omitted: $\binom42=6>1$ at $\rho=0,n=4$.

Let $B$ initially denote the proposed common CRT product and choose
\begin{equation}\label{eq:radix}
 q=B(V_H+1).
\end{equation}
The base-$q$ modulus row excludes $q=0$, so $B\ge1$ and $q>V_H$.
Lemma~\ref{lem:binomial} bounds all required lower coefficients of
$(1+q)^{\rho+H}$ by $V_H$, and those of $(1+q)^{\rho+t}$ by
$V_t\le V_H$, since $H\ge t$. The usual lower-digit extraction is
therefore exact for both factorial binomials before either factorial
has been interpreted.

For precision, if the coefficients through index $k$ are below $q$, then
the remainder of $(q+1)^n$ below $q^k$ lies in $[0,q^k)$ and
\[
 \floor{(q+1)^n/q^k}\equiv\binom nk\pmod q.
\]
The graph enforces the actual quotient, remainder and digit bounds by
quadratic equations. Coefficients above $k$ need no bound. If $B\mid q$,
divisibility of the quotient by $B$ is equivalent to
$B\mid\binom nk$; this is the fused range graph used later.

\section{The shared factorial bootstrap and the remaining radix bounds}
\begin{lemma}[Squared-output factorial identity]\label{lem:factorial}
For all $n,F\in\N$,
\[
 F=n!\quad\longleftrightarrow\quad
 F=\ceil{\frac{(F^2+1)^n}{\binom{F^2+n}{n}}}.
\]
\end{lemma}
\begin{proof}
Write $K=n!$ and
$f=K\prod_{j=1}^n(F^2+1)/(F^2+j)$. A fixed point has $1\le F\le K$.
For $n\ge2$, the finite-product inequality gives the strict deficit bound
\begin{equation}\label{eq:deficit}
 0\le K-f<\frac{K n(n-1)}{2F^2}.
\end{equation}
Indeed the sum of the factor deficiencies is strictly smaller than
$n(n-1)/(2F^2)$, and $\prod(1-a_j)\ge1-\sum a_j$.

For $n\ge4$, monotonicity in $F^2$ yields $f\ge2^n/(n+1)$.
At $n=4$ this forces $F\ge4$. For $n\ge5$ the elementary inequality
$2^n\ge n(n+1)$ forces $F\ge n$. Thus the relative deficit in
\eqref{eq:deficit} is less than one half, so $f>K/2$ and $F>K/2$.
Applying \eqref{eq:deficit} again gives
\[
 K-f<\frac{2n(n-1)}K\le1,
\]
because $n!\ge2n(n-1)$ for $n\ge4$. Therefore $\lceil f\rceil=K$.
For $n=0,1$, $f=1$. For $n=2$, $f=2(F^2+1)/(F^2+2)$ lies strictly
between one and two. For $n=3$ only $F=1,\ldots,6$ are possible; the
corresponding ceilings are $2,4,5,6,6,6$, so only six is fixed.

Conversely set $F=K$. For $n\ge2$, the deficit is less than
$n(n-1)/(2K)\le1/2$, and hence the ceiling is $K$. The cases $n=0,1$
were already handled.
\end{proof}
This theorem is inherited from the Bootstrap arrival; its proof is
repeated to make the present arithmetic argument self-contained.

Apply it first with $n=H,F=T$ after the exact binomial recovery above.
It gives $T=H!$. For $1\le t\le H$, the same common radius $\rho=(H!)^2$
has
\[
 0\le t!-\frac{(\rho+1)^t}{\binom{\rho+t}{t}}
 <\frac{t!\,t(t-1)}{2(H!)^2}
 \le\frac{H(H-1)}{2H!}\le\frac12
\]
when $t\ge2$. Therefore the second ceiling is $J=t!$.
For a correct ceiling $F=\lceil V/C\rceil$, completeness chooses
$r=V+C-CF-1$ and $s=C-r-1$, both nonnegative.

\subsection{The other binomials fit without assuming the CRT product}
Put $Q=(1+(t+1)T)^t+2$. The literal remainder rows already imply
$B<TQ$ before $B$ is known to be the intended product. The aggregate
range row bounds each of the eight positive column representatives
$a_i$ by $8B$. The unchanged elementary estimates give
\[
 Q+t,\ a_i<H^{(t+2)(H+1)}.
\]
For example, $T=H!\le H^H$ and $t+1\le H$ bound
$1+(t+1)T$ by $H^{H+2}$; the remaining factors $T,8$ and additive
margins are absorbed by the displayed larger exponent because
$H\ge R^4$, $R\ge4$ and $t\le R$. These are the same inequalities
proved in the pinned parent, evaluated at the new $H$.

Every needed lower index $k$ is at most $R$. Also
\[
 (H!)^2=\prod_{i=1}^{H}i(H+1-i)\ge H^H,
\]
since $i(H+1-i)-H=(i-1)(H-i)\ge0$. It follows that
\[
 V_H=((H!)^2+1)^H>H^{H^2}.
\]
From $R\ge4$ and $H\ge R^4$ we have $4R(R+2)\le H$; hence
\[
 k(t+2)(H+1)\le R(R+2)(H+1)<H^2.
\]
Each remaining lower coefficient is at most its upper argument to the
power $k$, and so is strictly below $V_H<q$. The product binomial and
all eight fused range graphs are now exact. The order matters: the
bound $B<TQ$ is a remainder consequence, not a use of the intended
CRT-product meaning of $B$.

\section{Closing the CRT and refutation equivalence}
After the arithmetic interpretations above, the moduli are
\[
 D_i=1+(i+1)H!\quad(0\le i<t),\qquad B=\prod_{i<t}D_i.
\]
They are pairwise coprime. A common divisor of two would divide their
index difference, which divides $H!$, whereas every $D_i$ is coprime
to $H!$.

The product is represented by the finite identity
\[
 B=\bigl(T^t t!\binom{Q+t}{t}\bigr)\bmod(TQ-1).
\]
Indeed $TQ\equiv1\pmod{TQ-1}$ turns the numerator into
$\prod_{j=1}^t(1+jT)$, and this true product is less than $TQ-1$ by
the definition of $Q$. All remainders and their bounds are actual rows.
The canonical encoded index is
$C=(B-1)/T-1\ge0$; reducing modulo $D_i$ gives $C\equiv i$.

For a column bound $b_j$, the range certificate requires
$B\mid\binom{a_j}{b_j+1}$. Choose any prime divisor $p$ of $D_i$.
Since $p$ is coprime to $H!$, $p>H$, in particular it exceeds every
column bound plus one. Multiplication by $(b_j+1)!$ shows that one
of the integers $a_j,a_j-1,\ldots,a_j-b_j$ is divisible by $p$.
Thus $a_j$ has a residue in $[0,b_j]$. The encoded local row is zero
modulo $B$, hence modulo $p$. Its decoded value is a nonnegative
integer strictly below $H<p$, by the local majorant, and therefore
is actually zero. This recovers every genuine local row at every
$i<t$.

The decoded accumulator coordinates satisfy $S_i,S'_i\le y$.
The modulus $1+(i+1)y$ is strictly larger than $y$, so its remainder is
unique. The second beta equation at index $i$ and the first beta equation
at index $i+1$ therefore force $S'_i=S_{i+1}$. The same uniqueness
argument joins the first and last rows to the endpoint equations.
The endpoint beta rows thus give one accumulator path, the carrier
lemma gives flow, and the flow lemma gives C222 refutability.

Conversely, from a finite refutation choose a padded $t\ge2$, then $e,M$
and the accumulator trace. Choose a suitable beta modulus $y$ and code
$x$, then $R,H$. Set $T=H!$ and $J=t!$, and choose the true product
$B$. CRT supplies positive canonical column representatives in $[1,B]$,
so their sum is at most $8B$. Their range divisibilities follow because
all relevant small factorials are invertible modulo each $D_i$.
Choose $C$, the shared radius, the radix, all true powers and all
Euclidean witnesses. Only then form $W$ from these values and choose
the common Pell frame and its finite-unit witnesses. Nothing at this
last stage changes an earlier arithmetic argument.

This gives both directions of
\[
 \exists\text{ natural raw arithmetic certificate}
 \quad\longleftrightarrow\quad\Refut(\mathrm{ZFC}).
\]
Every step uses finite arithmetic, induction, finite proof transformations
or the explicitly named Pell/CRT/calculus results. These are ordinary
mathematics formalizable in ZFC. Negating the equivalence gives
\eqref{eq:zfc}; it does not prove either side separately.

\section{From the certificate to the literal quartic}
The new raw graph has 1,232 natural coordinates and 1,149 quadratic
residuals. The exporter first substitutes 599 monic affine definitions
whose restoring polynomials have nonnegative integer coefficients.
It leaves 633 coordinates, 550 rows and a 3,552-term quartic.

A second exact contraction eliminates 275 more coordinates. Every
deleted coordinate has a monic defining row and an integer restoring
polynomial $f$. A certificate writes $f$ as a sum of positively weighted
rational squares of affine rational forms, plus a polynomial with
nonnegative integer coefficients. Thus $f\ge0$ on every natural tuple.
Its integer coefficients ensure that its values remain integral.
Two of the final restoration formulas have negative collected
coefficients; their explicit rational-square certificates are essential.

All rows are substituted at each step, and every remaining residual
is checked to have degree at most two. Reading the elimination list
backwards restores a natural solution of the parent. Reading it forwards
projects a parent solution to the child. No variable is silently
discarded as an unconstrained coordinate. This proves equivalence of
natural solution existence, not just agreement on sample points.

The final polynomial is exactly the sum of squares in \eqref{eq:main}.
For a simple exact-degree witness, the literal residual
$r_3=X_{18}-X_{17}^2$ contributes $X_{17}^4$; in a sum of squares,
the collected coefficient of that pure fourth power is a sum of squares
of the quadratic coefficients and is positive. Complete expansion also
finds and records every quartic coefficient. All 358 variables occur.

The 361-variable companion stops under a smaller support cap and has
6,759 monomials. The choice between these is a variable/support tradeoff,
not a difference in the consistency theorem.

\subsection{An explicit unrestricted-integer companion}
If an unrestricted integer domain is required, define
\[
 g(a,b,c)=a^2+b^2+c(c+1).
\]
It takes only nonnegative integer values and represents every natural
number. For surjectivity, apply the three-square theorem to $4n+1$:
exactly one square is odd, giving the displayed representation after
division by four. Consequently the explicit polynomial
\begin{equation}\label{eq:integer}
 \widehat P(\boldsymbol a,\boldsymbol b,\boldsymbol c)
 =\Pmain\bigl(g(a_0,b_0,c_0),\ldots,g(a_{357},b_{357},c_{357})\bigr)
\end{equation}
has 1,074 unrestricted integer variables and exact degree eight, and
ZFC proves its unsolvability equivalent to $\Con(\mathrm{ZFC})$.
The exact degree follows by setting all $b_i,c_i$ to zero: the nonzero
quartic part of $\Pmain$ becomes degree eight after $X_i=a_i^2$.
This companion is supplied as an explicit finite polynomial-composition
file and evaluator. It is a domain conversion, not an improvement over
the repository's separate integer quartics.

\section{Reproducibility, verification and limits}
The package includes the complete mathematical source, the two literal
quartics, their collected coefficients, raw graph and semantic inventory,
every projection certificate, the pinned upstream engine, source/calibration
files, and independent ordinary reviews. It is self-contained for Python
reconstruction; no network access or third-party Python package is needed.

The main commands, from the package root, are
\begin{verbatim}
python3 -B code/construction.py --check
python3 -B code/verify.py
python3 -B -O code/verify.py
python3 -B code/check_arithmetic.py
pdflatex -interaction=nonstopmode -halt-on-error report.tex
pdflatex -interaction=nonstopmode -halt-on-error report.tex
\end{verbatim}
The full reconstruction compares every emitted byte with the delivered
direct artifacts. The independent verifier imports no generator or
optimizer. It replays the entire raw-to-affine inverse, the majorant
weights, every rational-square restoration, every final row, all
collected coefficients and liveness. It also rejects an altered final
row even after the corresponding target hash is updated.

A separately written raw checker expands the shared-exponent identities,
the entire frame atom sum, the extra seed, its actual factor in the
auxiliary product, and all three addition-law pairs. The five evaluator
contractions have another independent checker which proves complete
parent/child output identities by local coefficient collection and
structural induction. Normal and optimized interpreter receipts are
retained separately.

The arithmetic controls additionally check 137,216 small binomial
instances, 5,663 factorial fixed-point candidates, 135 common-radius
remainder certificates and 111,881 radix exponent comparisons. These
finite controls are transcription and boundary checks. The universal
statements are justified by the proofs above, not by extrapolating tests.

The new equation has not been proved in Lean, Rocq or a sealed
object-language ZFC proof file. The inherited formal interfaces are
included for provenance, not represented as a rebuilt transitive formal
development. The result is an explicit equation with an ordinary proof
formalizable in ZFC, including its named inherited calibration. Neither
global simplicity nor consistency itself is established.

\section{Supplement: exact contractions of five existing equations}

This supplement gives five smaller arithmetic evaluators. Each computes
\emph{exactly the same polynomial} as its pinned parent, on every input and
for arbitrary values of the four coefficient ports. Thus the improvements
preserve the full zero sets, coefficient height, expanded support and exact
total degree. Their consistency interpretation is inherited from the
ordinary parent proofs; this supplement does not supply a new native
compiler or an internal proof-assistant formalization.

All unknowns are strictly positive integers. Binary addition, subtraction,
multiplication, squaring and multiplication by a scalar each cost one
operation; fixed constants and reuse cost zero.
\begin{center}\small
\begin{tabular}{rrrrrr}
Variables & Parent & New & Products & Sums/differences & Degree\\\hline
10 & 99 & 98 & 58 & $22+18$ & 368\\
10 & 105 & 103 & 61 & $23+19$ & 340\\
11 & 111 & 107 & 59 & $31+17$ & 211\\
12 & 96 & 95 & 52 & $25+18$ & 70\\
13 & 88 & 87 & 47 & $22+18$ & 51
\end{tabular}
\end{center}
The parents are the literal JSON equations at ProveIt commit
\file{84e16e43092614bae262ca59651cb9d77e31a255}; exact paths and SHA256
hashes are in \file{supplement/SOURCE_PINS.json}. Global optimality is not claimed.
The incoming Bootstrap package already evaluates this eleven-variable
polynomial in 110 operations at degree 211; the new 107-operation circuit
saves three relative to that point and four relative to its literal
111-operation parent.

\subsection{Contraction of an affine carrier}
For arbitrary elements of any commutative ring,
\begin{equation}
 pqz+qY-p(1+qL)B=pq(z-LB)+qY-pB.
 \label{eq:supp-affine}
\end{equation}
This follows by distributing both products. When $pq$ is already used in
another part of the equation, the right side avoids constructing the
carrier $1+qL$: four old operations are replaced by three.

The scaled ten-variable version is the equally exact identity
\begin{align*}
 &pqu(Hz+K)+H\{u[qY-p(1+qL)B]-pqe\}\\
 &\qquad=pqu\{H(z-LB)+K\}+H\{u(qY-pB)-pqe\}.
\end{align*}
In the degree-368 equation set $p=q_{\rm source}^3$,
$q=q_{\rm source}-1$, $u=\mu^2$, $z=n+V$, $B=B_t$ and $K=K_{\rm pos}$.
No unknown is shifted or reparametrized: $z-LB$ is only an internal
polynomial abbreviation. It need not be positive. This distinction means
that no additional unbounded witness argument is required.

For the degree-340 direct form the carrier coefficient is
$q_{\rm source}(q_{\rm source}-1)$ and one further exact factorization is
used:
\[
 q_{\rm source}(q_{\rm source}-1)Y-(q_{\rm source}-1)X
  =(q_{\rm source}-1)(q_{\rm source}Y-X).
\]
After contraction the former intermediate
$q_{\rm source}(q_{\rm source}-1)$ has no other consumer. Eliminating it
saves the second operation. The degree-70 and degree-51 equations use
\eqref{eq:supp-affine} directly in their first Pell carrier.

\subsection{A shorter norm in the eleven-variable evaluator}
Define
\[
 \mathcal N(s,a,b,c)
 =\bigl((s+c-a-b)^2-4(sc+ab)\bigr)^2-64scab.
\]
Sharing $sc$ and $ab$ gives a 13-operation evaluator. For arbitrary
$z,a,b,c,\delta$ the following identity holds:
\begin{align*}
 &\left(\bigl[\delta(z-a-b)-c\bigr]^2
        -4[\delta^2ab+\delta c(a+b)]\right)^2
       -64\delta^3abc z\\
 &\qquad=\mathcal N(\delta z,\delta a,\delta b,c).
 \label{eq:supp-norm}
\end{align*}
To check it, put $J=\delta(z-a-b)-c$. The inner expression on the right is
\[
 (J+2c)^2-4(\delta zc+\delta^2ab)
 =J^2-4[\delta^2ab+\delta c(a+b)],
\]
and the final product is the same $64\delta^3abc z$.
Thus the old 18-operation tail after $z=\overline Z^2$ becomes three
scalings followed by the 13-operation norm, saving two operations.
One more operation is saved by reusing $E=XY$ in
$U=q^2Y$ and $qE$, instead of first constructing the private intermediate
$S=qY$. Combined with the affine-carrier contraction, this gives
$111-2-1-1=107$ operations. The eleven-variable restriction
$J=3j+1$ is retained exactly; no enlarged-domain theorem is asserted.

\subsection{Exact and reproducible verification}
The package supplies all five new circuits, the five complete pinned
parent circuits, all four actual coefficient payloads, their reconstruction
recipe, inherited proof notes and source notices. The coefficient integers
are represented without omission as unsigned minimal little-endian bytes
inside gzip files. Their individual bit lengths are 590069, 592884, 590069
and 590070 for $\beta,K_0,M_C,M_F$, respectively. In particular
$\beta=2^{590069}-1$. The already shifted $M_F$ is used unchanged.

The independent \texttt{check.py} imports neither builder. It reconstructs
small free-polynomial identities directly from the literal parent and child
gates, collects every coefficient over the integers, and then propagates
certified register equality through the entire circuit. The five complete
output equalities follow by structural induction. It also verifies
acyclicity, complete liveness, exact operation counts, every source hash,
every compressed and raw coefficient hash, and rejection of altered output
gates. The largest local comparison has 35 collected terms. Normal and
optimized Python runs both pass. Modular controls in the builders are
supplementary diagnostics and are not the identity proof.

Exact degrees transfer from the pinned parents by the proved polynomial
identities. The larger syntax-only bounds, respectively 416, 384, 262, 82
and 58, are deliberately distinguished from the exact degrees in the table.
This supplement establishes a finite algebraic optimization, not a proof
that any of these polynomials lacks zeros.


\section{What the reduction suggests next}
The successful optimization treats the arithmetic quantities as
certificates with freedom of choice. The large bound $H$ only needed
to dominate a source; requiring its original exact value paid for an
unnecessary index. Likewise, the pair at $\rho$ was useful even though
no arithmetic request asked for that exponent alone. These observations
suggest optimizing the finite set of Pell indices as an addition graph,
allowing intermediate indices, before minimizing its polynomial wires.

A next implementation should search such index graphs while preserving
the positivity and dependency order proved here. It should also keep
coefficient height in its objective: absorbing fixed exponents into
scalar coefficients saved variables but raised height from 27 to 94
bits. A better point may retain selected squaring wires while sharing
the remaining indices. This is a concrete next target, not a claimed
further reduction in the present release.

\begin{thebibliography}{9}
\bibitem{repo} V. Reshetnikov and ProveIt contributors.
\emph{ProveIt}, snapshot
\file{84e16e43092614bae262ca59651cb9d77e31a255}.
\url{https://github.com/VladimirReshetnikov/ProveIt/tree/84e16e43092614bae262ca59651cb9d77e31a255}.

\bibitem{frontier} ProveIt, \emph{Later arrivals and current size reductions},
\file{Computability/DiophantineConsistency/PROPOSAL_REVIEW_LATEST.md},
at the same snapshot. A local copy is \file{proofs/pinned_frontier.md}.

\bibitem{directreview} ProveIt, \emph{Ordinary correctness review: direct C222
quartic, round 12},
\file{Computability/DiophantineConsistency/PROPOSAL_REVIEW_DIRECT_C222_ROUND12.md}.
The local copy is \file{proofs/parent_direct_calibration_review.md}.

\bibitem{parent407} \emph{ConZFC\_407\_888\_Quartics\_2026-10-24.zip},
in the pinned repository's \file{docs/incoming/}. SHA256:
\file{59f847872298260ee35e9d8e4cd0b52d2d78658ac523981fe060e81f40ca4a59}.
Included under \file{upstream/engine/}; original proofs and source data retained.

\bibitem{parent399} ProveIt, \emph{Projected low-digit quartics with
candidate-output factorial radii},
\file{Computability/DiophantineConsistency/candidates/projected-factorial-radius/}.
The parent proof is copied to \file{proofs/parent_projected_factorial_radius.md}.

\bibitem{bootstrap} \emph{ConZFC\_Bootstrap\_Certificates\_2026-10-24.zip},
\file{docs/incoming/}, especially \file{direct/proofs/NEW_RESULTS.md},
Sections 2--3. The squared-output factorial theorem and shared radius
predate this report.

\bibitem{centered} \emph{ConZFC\_Valuations\_and\_Centered\_Certificates.zip},
\file{docs/incoming/}, especially the centered factorial proof and
395-variable natural quartic. Comparison scope is recorded in
\file{proofs/incoming_comparison.md}.
\end{thebibliography}

\appendix
\clearpage
\section{The complete 275-residual definition}\label{app:residuals}
All variables below range over $\N$. Every occurrence of $X_i$ uses the
zero-based coordinate order of the delivered JSON. These are the entire
polynomials $r_0,\ldots,r_{274}$, including all constants and signs;
the equation is their sum of squares equal to zero as in \eqref{eq:main}.
They were transcribed automatically from the verified residual list.
The separately stored 8,229-term expansion is obtained by collecting
\emph{all} products in that sum of squares.

% Generated by code/render_residuals.py from direct/natural358/equation.json.
% Complete integer coefficients; no omitted terms or symbolic constants.
\begingroup\small
\setlength{\jot}{2pt}
\allowdisplaybreaks[4]
\begin{align*}
r_{0} &= -6 - 2X_{0} - X_{0}X_{2} - X_{0}X_{4} - 3X_{2} - 3X_{4} + X_{6} \\[3pt]
r_{1} &= -2 + X_{1} - X_{2} - X_{2}X_{3} - 3X_{3} - X_{3}X_{4} \\[3pt]
r_{2} &= -1 + X_{1} - X_{5} - X_{5}X_{6} \\[3pt]
r_{3} &= -X_{17}^{2} + X_{18} \\[3pt]
r_{4} &= -X_{17} - X_{17}X_{20} - X_{17}X_{21} + X_{19} + X_{21} \\[3pt]
r_{5} &= -X_{0}X_{17} - 3X_{17} + X_{22} \\[3pt]
r_{6} &= -X_{25} - X_{25}X_{27} - X_{25}X_{28} + X_{26} + X_{28} \\[3pt]
r_{7} &= -X_{24}X_{25} + X_{30} \\[3pt]
r_{8} &= 2 - 2X_{17} - X_{17}X_{23} + X_{31} + X_{33} \\[3pt]
r_{9} &= X_{29}X_{30} - X_{31} - X_{31}X_{32} - X_{32} - X_{32}X_{33} \\[3pt]
r_{10} &= 1 + X_{15}X_{17} + X_{17} - X_{31} \\[3pt]
r_{11} &= -6 - 5X_{7} - X_{7}^{2} - 2X_{7}X_{8} - 5X_{8} - X_{8}^{2} + X_{34} \\[3pt]
r_{12} &= -72 - 68X_{8} - 16X_{8}^{2} - 16X_{8}X_{34} - 34X_{34} - 4X_{34}^{2} + X_{35} \\[3pt]
r_{13} &= -30 - 11X_{7} - X_{7}^{2} - 4X_{7}X_{35} - 22X_{35} - 4X_{35}^{2} + X_{36} \\[3pt]
r_{14} &= -380 - 312X_{35} - 64X_{35}^{2} - 32X_{35}X_{36} - 78X_{36} - 4X_{36}^{2} + X_{37} \\[3pt]
r_{15} &= -6 - 5X_{8} - X_{8}^{2} - 2X_{8}X_{9} - 5X_{9} - X_{9}^{2} + X_{38} \\[3pt]
r_{16} &= -6 - 5X_{7} - X_{7}^{2} - 2X_{7}X_{9} - 5X_{9} - X_{9}^{2} + X_{39} \\[3pt]
r_{17} &= -414 - 22X_{7} - 2X_{7}^{2} - 16X_{7}X_{9} - 8X_{7}X_{38} - 200X_{8} - 32X_{8}^{2} - 64X_{8}X_{9} - 32X_{8}X_{34} - 32X_{8}X_{39} \\*
&\quad  - 320X_{9} - 64X_{9}^{2} - 32X_{9}X_{34} - 32X_{9}X_{38} - 32X_{9}X_{39} - 100X_{34} - 8X_{34}^{2} - 16X_{34}X_{39} - 52X_{38} \\*
&\quad  - 8X_{38}^{2} - 108X_{39} - 8X_{39}^{2} + X_{40} \\[3pt]
r_{18} &= -1344 - 800X_{8} - 128X_{8}^{2} - 256X_{8}X_{9} - 128X_{8}X_{34} - 128X_{8}X_{39} - 864X_{9} - 128X_{9}^{2} \\*
&\quad  - 128X_{9}X_{34} - 128X_{9}X_{39} - 400X_{34} - 32X_{34}^{2} - 64X_{34}X_{39} - 432X_{39} - 32X_{39}^{2} - 2X_{40} - 2X_{40}^{2} \\*
&\quad  + X_{41} \\[3pt]
r_{19} &= -247 - 168X_{7} - 32X_{7}^{2} - 64X_{7}X_{8} - 32X_{7}X_{38} - 32X_{7}X_{39} - 188X_{8} - 32X_{8}^{2} - 32X_{8}X_{38} \\*
&\quad  - 32X_{8}X_{39} - 2X_{34} - 92X_{38} - 8X_{38}^{2} - 16X_{38}X_{39} - 84X_{39} - 8X_{39}^{2} + X_{42} \\[3pt]
r_{20} &= -960 - 672X_{7} - 128X_{7}^{2} - 256X_{7}X_{8} - 128X_{7}X_{38} - 128X_{7}X_{39} - 736X_{8} - 128X_{8}^{2} - 128X_{8}X_{38} \\*
&\quad  - 128X_{8}X_{39} - 368X_{38} - 32X_{38}^{2} - 64X_{38}X_{39} - 336X_{39} - 32X_{39}^{2} - 2X_{42} - 2X_{42}^{2} + X_{43} \\[3pt]
r_{21} &= -276 - 30X_{7} - 2X_{7}^{2} - 16X_{7}X_{9} - 8X_{7}X_{38} - 30X_{8} - 2X_{8}^{2} - 16X_{8}X_{9} - 8X_{8}X_{39} - 272X_{9} \\*
&\quad  - 64X_{9}^{2} - 32X_{9}X_{38} - 32X_{9}X_{39} - 68X_{38} - 8X_{38}^{2} - 68X_{39} - 8X_{39}^{2} + X_{44} \\[3pt]
r_{22} &= -552 - 120X_{8} - 8X_{8}^{2} - 64X_{8}X_{9} - 32X_{8}X_{39} - 544X_{9} - 128X_{9}^{2} - 128X_{9}X_{39} - 272X_{39} \\*
&\quad  - 32X_{39}^{2} - 2X_{44} - 2X_{44}^{2} + X_{45} \\[3pt]
r_{23} &= -247 - 4X_{7} - 352X_{8} - 128X_{8}^{2} - 64X_{8}X_{34} - 64X_{8}X_{38} - 92X_{34} - 8X_{34}^{2} - 16X_{34}X_{38} - 84X_{38} \\*
&\quad  - 8X_{38}^{2} - 2X_{39} + X_{46} \\[3pt]
r_{24} &= -960 - 1408X_{8} - 512X_{8}^{2} - 256X_{8}X_{34} - 256X_{8}X_{38} - 368X_{34} - 32X_{34}^{2} - 64X_{34}X_{38} - 336X_{38} \\*
&\quad  - 32X_{38}^{2} - 2X_{46} - 2X_{46}^{2} + X_{47} \\[3pt]
r_{25} &= -56 - 60X_{8} - 16X_{8}^{2} - 16X_{8}X_{34} - 30X_{34} - 4X_{34}^{2} + X_{48} \\[3pt]
r_{26} &= -30 - 11X_{7} - X_{7}^{2} - 4X_{7}X_{48} - 22X_{48} - 4X_{48}^{2} + X_{49} \\[3pt]
r_{27} &= -380 - 312X_{48} - 64X_{48}^{2} - 32X_{48}X_{49} - 78X_{49} - 4X_{49}^{2} + X_{50} \\[3pt]
r_{28} &= -30 - 11X_{8} - X_{8}^{2} - 4X_{8}X_{50} - 22X_{50} - 4X_{50}^{2} + X_{51} \\[3pt]
r_{29} &= -349 - 420X_{8} - 128X_{8}^{2} - 128X_{8}X_{34} - 208X_{34} - 32X_{34}^{2} - 2X_{38} + X_{52} \\[3pt]
r_{30} &= -1368 - 1664X_{8} - 512X_{8}^{2} - 512X_{8}X_{34} - 832X_{34} - 128X_{34}^{2} - 2X_{52} - 2X_{52}^{2} + X_{53} \\[3pt]
r_{31} &= -56 - 60X_{7} - 16X_{7}^{2} - 16X_{7}X_{39} - 8X_{7}X_{53} - 30X_{39} - 4X_{39}^{2} - 4X_{39}X_{53} - 15X_{53} - X_{53}^{2} \\*
&\quad  + X_{54} \\[3pt]
r_{32} &= -6 - 5X_{9} - X_{9}^{2} - 2X_{9}X_{10} - 5X_{10} - X_{10}^{2} + X_{55} \\[3pt]
r_{33} &= -224 - 26X_{7} - 2X_{7}^{2} - 16X_{7}X_{10} - 8X_{7}X_{55} - 288X_{9} - 128X_{9}^{2} - 64X_{9}X_{38} - 64X_{9}X_{39} - 88X_{10} \\*
&\quad  - 32X_{10}^{2} - 32X_{10}X_{55} - 76X_{38} - 8X_{38}^{2} - 16X_{38}X_{39} - 68X_{39} - 8X_{39}^{2} - 44X_{55} - 8X_{55}^{2} + X_{56} \\[3pt]
r_{34} &= -640 - 1152X_{9} - 512X_{9}^{2} - 256X_{9}X_{38} - 256X_{9}X_{39} - 304X_{38} - 32X_{38}^{2} - 64X_{38}X_{39} - 272X_{39} \\*
&\quad  - 32X_{39}^{2} - 2X_{56} - 2X_{56}^{2} + X_{57} \\[3pt]
r_{35} &= -241 - 168X_{7} - 32X_{7}^{2} - 64X_{7}X_{8} - 32X_{7}X_{38} - 32X_{7}X_{39} - 184X_{8} - 32X_{8}^{2} - 32X_{8}X_{38} \\*
&\quad  - 32X_{8}X_{39} - X_{9} - 92X_{38} - 8X_{38}^{2} - 16X_{38}X_{39} - 84X_{39} - 8X_{39}^{2} + X_{58} \\[3pt]
r_{36} &= -42 - 52X_{9} - 16X_{9}^{2} - 16X_{9}X_{38} - 26X_{38} - 4X_{38}^{2} + X_{59} \\[3pt]
r_{37} &= -30 - 11X_{9} - X_{9}^{2} - 4X_{9}X_{59} - 22X_{59} - 4X_{59}^{2} + X_{60} \\[3pt]
r_{38} &= -306 - 35X_{8} - X_{8}^{2} - 16X_{8}X_{59} - 4X_{8}X_{60} - 280X_{59} - 64X_{59}^{2} - 32X_{59}X_{60} - 70X_{60} - 4X_{60}^{2} \\*
&\quad  + X_{61} \\[3pt]
r_{39} &= -4556 - 4320X_{59} - 1024X_{59}^{2} - 512X_{59}X_{60} - 128X_{59}X_{61} - 1080X_{60} - 64X_{60}^{2} - 32X_{60}X_{61} \\*
&\quad  - 270X_{61} - 4X_{61}^{2} + X_{62} \\[3pt]
r_{40} &= -30 - 11X_{9} - X_{9}^{2} - 4X_{9}X_{62} - 22X_{62} - 4X_{62}^{2} + X_{63} \\[3pt]
r_{41} &= -306 - 35X_{8} - X_{8}^{2} - 16X_{8}X_{62} - 4X_{8}X_{63} - 280X_{62} - 64X_{62}^{2} - 32X_{62}X_{63} - 70X_{63} - 4X_{63}^{2} \\*
&\quad  + X_{64} \\[3pt]
r_{42} &= -1028 - 672X_{7} - 128X_{7}^{2} - 256X_{7}X_{8} - 128X_{7}X_{38} - 128X_{7}X_{39} - 736X_{8} - 128X_{8}^{2} \\*
&\quad  - 128X_{8}X_{38} - 128X_{8}X_{39} - 368X_{38} - 32X_{38}^{2} - 64X_{38}X_{39} - 336X_{39} - 32X_{39}^{2} - 2X_{58} - 2X_{58}^{2} \\*
&\quad  - 32X_{62} - 8X_{63} - 2X_{64} + X_{65} \\[3pt]
r_{43} &= -264 - 128X_{62} - 32X_{63} - 8X_{64} - 2X_{65} - 2X_{65}^{2} + X_{66} \\[3pt]
r_{44} &= -2 - 3X_{7} - X_{7}^{2} - 2X_{7}X_{66} - 3X_{66} - X_{66}^{2} + X_{67} \\[3pt]
r_{45} &= -12 - 28X_{66} - 16X_{66}^{2} - 16X_{66}X_{67} - 14X_{67} - 4X_{67}^{2} + X_{68} \\[3pt]
r_{46} &= -30 - 11X_{9} - X_{9}^{2} - 4X_{9}X_{68} - 22X_{68} - 4X_{68}^{2} + X_{69} \\[3pt]
r_{47} &= -306 - 35X_{7} - X_{7}^{2} - 16X_{7}X_{68} - 4X_{7}X_{69} - 280X_{68} - 64X_{68}^{2} - 32X_{68}X_{69} - 70X_{69} - 4X_{69}^{2} \\*
&\quad  + X_{70} \\[3pt]
r_{48} &= -4556 - 4320X_{68} - 1024X_{68}^{2} - 512X_{68}X_{69} - 128X_{68}X_{70} - 1080X_{69} - 64X_{69}^{2} - 32X_{69}X_{70} \\*
&\quad  - 270X_{70} - 4X_{70}^{2} + X_{71} \\[3pt]
r_{49} &= -90 - 76X_{9} - 16X_{9}^{2} - 16X_{9}X_{38} - 16X_{9}X_{71} - 38X_{38} - 4X_{38}^{2} - 8X_{38}X_{71} - 38X_{71} - 4X_{71}^{2} \\*
&\quad  + X_{72} \\[3pt]
r_{50} &= -4556 - 4320X_{62} - 1024X_{62}^{2} - 512X_{62}X_{63} - 128X_{62}X_{64} - 1080X_{63} - 64X_{63}^{2} - 32X_{63}X_{64} \\*
&\quad  - 270X_{64} - 4X_{64}^{2} + X_{73} \\[3pt]
r_{51} &= -682 - 70X_{8} - 2X_{8}^{2} - 32X_{8}X_{71} - 8X_{8}X_{72} - 592X_{71} - 128X_{71}^{2} - 64X_{71}X_{72} - 148X_{72} \\*
&\quad  - 8X_{72}^{2} - 2X_{73} + X_{74} \\[3pt]
r_{52} &= -X_{74} - X_{74}^{2} + X_{75} \\[3pt]
r_{53} &= -42 - 52X_{8} - 16X_{8}^{2} - 8X_{8}X_{9} - 16X_{8}X_{34} - 13X_{9} - X_{9}^{2} - 4X_{9}X_{34} - 26X_{34} - 4X_{34}^{2} + X_{76} \\[3pt]
r_{54} &= -2278 - 2264X_{8} - 800X_{8}^{2} - 800X_{8}X_{34} - 160X_{8}X_{76} - 70X_{9} - 2X_{9}^{2} - 32X_{9}X_{73} - 8X_{9}X_{75} \\*
&\quad  - 1132X_{34} - 200X_{34}^{2} - 80X_{34}X_{76} - 592X_{73} - 128X_{73}^{2} - 64X_{73}X_{75} - 148X_{75} - 8X_{75}^{2} \\*
&\quad  - 228X_{76} - 8X_{76}^{2} + X_{77} \\[3pt]
r_{55} &= -2712 - 280X_{9} - 8X_{9}^{2} - 128X_{9}X_{73} - 32X_{9}X_{75} - 2368X_{73} - 512X_{73}^{2} - 256X_{73}X_{75} - 592X_{75} \\*
&\quad  - 32X_{75}^{2} - 2X_{77} - 2X_{77}^{2} + X_{78} \\[3pt]
r_{56} &= -2 - 3X_{78} - X_{78}^{2} + X_{79} \\[3pt]
r_{57} &= -1 - X_{2} - X_{2}X_{12} - X_{12} + X_{80} \\[3pt]
r_{58} &= -X_{11}^{2} + 2X_{11}X_{12} - X_{12}^{2} + X_{81} \\[3pt]
r_{59} &= -2X_{37}X_{80} - 4X_{80} + X_{82} \\[3pt]
r_{60} &= -X_{11}^{2} + 2X_{11}X_{12} + 2X_{11}X_{82} - X_{12}^{2} - 2X_{12}X_{82} - X_{82}^{2} + X_{83} \\[3pt]
r_{61} &= -X_{41}X_{80} + X_{84} \\[3pt]
r_{62} &= -X_{11}^{2} + 2X_{11}X_{12} + 2X_{11}X_{84} - X_{12}^{2} - 2X_{12}X_{84} - X_{84}^{2} + X_{85} \\[3pt]
r_{63} &= -X_{43}X_{80} + X_{86} \\[3pt]
r_{64} &= -X_{11}^{2} + 2X_{11}X_{12} + 2X_{11}X_{86} - X_{12}^{2} - 2X_{12}X_{86} - X_{86}^{2} + X_{87} \\[3pt]
r_{65} &= -X_{45}X_{80} + X_{88} \\[3pt]
r_{66} &= -X_{11}^{2} + 2X_{11}X_{12} + 2X_{11}X_{88} - X_{12}^{2} - 2X_{12}X_{88} - X_{88}^{2} + X_{89} \\[3pt]
r_{67} &= -X_{47}X_{80} + X_{90} \\[3pt]
r_{68} &= -X_{11}^{2} + 2X_{11}X_{12} + 2X_{11}X_{90} - X_{12}^{2} - 2X_{12}X_{90} - X_{90}^{2} + X_{91} \\[3pt]
r_{69} &= -4X_{8}X_{80} - 2X_{51}X_{80} - 4X_{80} + X_{92} \\[3pt]
r_{70} &= -X_{11}^{2} + 2X_{11}X_{12} + 2X_{11}X_{92} - X_{12}^{2} - 2X_{12}X_{92} - X_{92}^{2} + X_{93} \\[3pt]
r_{71} &= -4X_{53}X_{80} - 2X_{54}X_{80} + X_{94} \\[3pt]
r_{72} &= -X_{11}^{2} + 2X_{11}X_{12} + 2X_{11}X_{94} - X_{12}^{2} - 2X_{12}X_{94} - X_{94}^{2} + X_{95} \\[3pt]
r_{73} &= -X_{57}X_{80} + X_{96} \\[3pt]
r_{74} &= -X_{11}^{2} + 2X_{11}X_{12} + 2X_{11}X_{96} - X_{12}^{2} - 2X_{12}X_{96} - X_{96}^{2} + X_{97} \\[3pt]
r_{75} &= -X_{7}^{2} + 2X_{7}X_{9} - X_{9}^{2} + X_{98} \\[3pt]
r_{76} &= -X_{8}^{2} + 2X_{8}X_{9} - X_{9}^{2} + X_{99} \\[3pt]
r_{77} &= -X_{98}X_{99} + X_{100} \\[3pt]
r_{78} &= -4 - 4X_{10} - X_{10}^{2} + X_{101} \\[3pt]
r_{79} &= -X_{100}^{2} + 2X_{100}X_{101} - X_{101}^{2} + X_{102} \\[3pt]
r_{80} &= -2X_{79}X_{80} - 4X_{80} + X_{103} \\[3pt]
r_{81} &= -X_{11}^{2} + 2X_{11}X_{12} + 2X_{11}X_{103} - X_{12}^{2} - 2X_{12}X_{103} - X_{103}^{2} + X_{104} \\[3pt]
r_{82} &= -1 - X_{2} - X_{2}X_{8} - X_{8} + X_{105} \\[3pt]
r_{83} &= -1 - X_{12} - X_{12}X_{105} - X_{105} + X_{106} \\[3pt]
r_{84} &= -1 - X_{2} - X_{2}X_{7} - X_{7} + X_{107} \\[3pt]
r_{85} &= -4 - 4X_{2} - 4X_{2}X_{8} - 2X_{2}X_{34} - 4X_{8} - 2X_{34} + X_{108} \\[3pt]
r_{86} &= -1 - X_{11} - X_{11}X_{107} - X_{107} + X_{109} \\[3pt]
r_{87} &= -X_{108}X_{109} - X_{109} + X_{110} \\[3pt]
r_{88} &= -X_{106}^{2} + 2X_{106}X_{110} - X_{110}^{2} + X_{111} \\[3pt]
r_{89} &= -6 - 6X_{2} - 4X_{2}X_{8} - 2X_{2}X_{34} - 4X_{8} - 2X_{34} + X_{112} \\[3pt]
r_{90} &= -1 - X_{12} - X_{12}X_{112} - X_{112} + X_{113} \\[3pt]
r_{91} &= -1 - X_{11} - X_{11}X_{105} - X_{105} + X_{114} \\[3pt]
r_{92} &= -X_{113}^{2} + 2X_{113}X_{114} - X_{114}^{2} + X_{115} \\[3pt]
r_{93} &= -X_{81}X_{83} + X_{116} \\[3pt]
r_{94} &= -X_{85}X_{116} + X_{117} \\[3pt]
r_{95} &= -X_{87}X_{117} + X_{118} \\[3pt]
r_{96} &= -X_{89}X_{118} + X_{119} \\[3pt]
r_{97} &= -X_{91}X_{119} + X_{120} \\[3pt]
r_{98} &= -X_{93}X_{120} + X_{121} \\[3pt]
r_{99} &= -X_{95}X_{121} + X_{122} \\[3pt]
r_{100} &= -X_{97}X_{122} + X_{123} \\[3pt]
r_{101} &= -X_{102}X_{123} - X_{104}X_{123} + X_{124} \\[3pt]
r_{102} &= -X_{111}X_{124} + X_{125} \\[3pt]
r_{103} &= -2 - X_{2} - X_{2}X_{15} - X_{4} - X_{4}X_{15} - 2X_{15} + X_{126} \\[3pt]
r_{104} &= -1 - X_{13} - X_{13}X_{126} - X_{126} + X_{127} \\[3pt]
r_{105} &= -3 - X_{2} - X_{2}X_{14} - X_{4} - X_{4}X_{14} - 3X_{14} - X_{14}X_{126} - X_{126} + X_{128} \\[3pt]
r_{106} &= 2 - 4X_{1} + 2X_{1}^{2} - 2X_{1}X_{11} - 2X_{1}X_{12} - 2X_{1}X_{127} - 2X_{1}X_{128} + 2X_{11} + X_{11}^{2} + 2X_{11}X_{127} \\*
&\quad  + 2X_{12} + X_{12}^{2} + 2X_{12}X_{128} - X_{31}X_{129} + X_{115}X_{125} + 2X_{127} + X_{127}^{2} + 2X_{128} + X_{128}^{2} \\[3pt]
r_{107} &= -8 - X_{7} - X_{8} - X_{9} - X_{10} - X_{11} - X_{12} - X_{13} - X_{14} + 8X_{31} - X_{130} \\[3pt]
r_{108} &= 2 - X_{19}X_{31} + X_{20} + X_{21} - X_{31} + X_{134} \\[3pt]
r_{109} &= -X_{20}X_{131} - X_{21}X_{131} - 2X_{131} - X_{131}X_{134} + X_{135} \\[3pt]
r_{110} &= -X_{20}X_{136} - X_{21}X_{136} - X_{134}X_{136} - 2X_{136} + X_{140} \\[3pt]
r_{111} &= X_{20} + X_{21} - X_{27} - X_{28} + X_{134} - X_{139} \\[3pt]
r_{112} &= -X_{20}X_{141} - X_{21}X_{141} - X_{134}X_{141} - 2X_{141} + X_{145} \\[3pt]
r_{113} &= X_{137} + X_{138} - X_{142} - X_{143} \\[3pt]
r_{114} &= 1 + X_{20} + X_{21} - X_{29} + X_{134} - X_{144} \\[3pt]
r_{115} &= -X_{31} - X_{31}X_{147} - X_{31}X_{148} + X_{149} \\[3pt]
r_{116} &= X_{147} + X_{148} - X_{151} - X_{152} \\[3pt]
r_{117} &= X_{147} + X_{148} - X_{154} - X_{155} \\[3pt]
r_{118} &= X_{147} + X_{148} - X_{157} - X_{158} \\[3pt]
r_{119} &= -X_{31} - X_{31}X_{160} - X_{31}X_{161} + X_{162} \\[3pt]
r_{120} &= X_{160} + X_{161} - X_{164} - X_{165} \\[3pt]
r_{121} &= -X_{31} - X_{31}X_{167} - X_{31}X_{168} + X_{169} \\[3pt]
r_{122} &= X_{167} + X_{168} - X_{171} - X_{172} \\[3pt]
r_{123} &= -X_{173} - X_{173}X_{175} + X_{176} \\[3pt]
r_{124} &= -X_{176}^{2} + X_{177} \\[3pt]
r_{125} &= 1 - X_{174}^{2} + X_{230} \\[3pt]
r_{126} &= -1 - 2X_{173} - X_{173}^{2} + X_{178} \\[3pt]
r_{127} &= X_{177} - X_{177}X_{178} + X_{230} \\[3pt]
r_{128} &= -X_{180}^{2} + X_{183} \\[3pt]
r_{129} &= -1 + X_{179}^{2} - X_{183}X_{230} \\[3pt]
r_{130} &= -4X_{180} - 4X_{180}X_{181} + X_{184} \\[3pt]
r_{131} &= 1 - X_{174} - X_{179}X_{182} + X_{184} \\[3pt]
r_{132} &= -70368744177664 - 70368744177664X_{0} - 17592186044416X_{0}^{2} - 8388608X_{0}X_{188} \\*
&\quad  - 16777216X_{188} - X_{188}^{2} + X_{189} \\[3pt]
r_{133} &= -1 + X_{185}^{2} - X_{189}X_{230} \\[3pt]
r_{134} &= -X_{187}^{2} + X_{190} \\[3pt]
r_{135} &= -1 - 2X_{184} - X_{184}^{2} + X_{191} \\[3pt]
r_{136} &= -1 + X_{186}^{2} + X_{190} - X_{190}X_{191} \\[3pt]
r_{137} &= -X_{179}X_{192} - X_{185} + X_{186} \\[3pt]
r_{138} &= -8388608 - 4194304X_{0} - 16777216X_{0}X_{193} + X_{187} - 4X_{188}X_{193} - 33554432X_{193} \\[3pt]
r_{139} &= -10 - 6X_{0} - X_{0}^{2} + 2X_{0}X_{174} + 6X_{174} - X_{194} \\[3pt]
r_{140} &= 25165823 + 20971520X_{0} + 4194304X_{0}^{2} - 4194304X_{0}X_{174} + X_{0}X_{188} - X_{2} - 8388608X_{174} \\*
&\quad  - X_{174}X_{188} + X_{185} + 3X_{188} - X_{194}X_{195} \\[3pt]
r_{141} &= -37 - 12X_{0} - X_{0}^{2} - 2X_{0}X_{1} - 4X_{0}X_{2} - 2X_{0}X_{4} + 2X_{0}X_{174} - 12X_{1} - X_{1}^{2} - 4X_{1}X_{2} - 2X_{1}X_{4} \\*
&\quad  + 2X_{1}X_{174} - 24X_{2} - 4X_{2}^{2} - 4X_{2}X_{4} + 4X_{2}X_{174} - 12X_{4} - X_{4}^{2} + 2X_{4}X_{174} + 12X_{174} - X_{196} \\[3pt]
r_{142} &= 50331648 + 33554432X_{0} + 4194304X_{0}^{2} + 4194304X_{0}X_{1} + 8388608X_{0}X_{2} + 4194304X_{0}X_{4} \\*
&\quad  - 4194304X_{0}X_{174} + X_{0}X_{188} + 8388608X_{1} + X_{1}X_{188} + 16777216X_{2} + 2X_{2}X_{188} + 8388608X_{4} \\*
&\quad  + X_{4}X_{188} - X_{16} - 8388608X_{174} - X_{174}X_{188} + X_{185} + 6X_{188} - X_{196}X_{197} \\[3pt]
r_{143} &= -X_{16}^{2} - 2X_{16}X_{201} - X_{201}^{2} + X_{202} \\[3pt]
r_{144} &= -1 + X_{198}^{2} - X_{202}X_{230} \\[3pt]
r_{145} &= -X_{200}^{2} + X_{203} \\[3pt]
r_{146} &= -1 - X_{191}X_{203} + X_{199}^{2} + X_{203} \\[3pt]
r_{147} &= -X_{179}X_{204} - X_{198} + X_{199} \\[3pt]
r_{148} &= -X_{16} - 4X_{16}X_{205} + X_{200} - 4X_{201}X_{205} \\[3pt]
r_{149} &= -2 - 2X_{18} - X_{18}^{2} + 2X_{18}X_{174} + 2X_{174} - X_{206} \\[3pt]
r_{150} &= X_{16} + X_{16}X_{18} - X_{16}X_{174} + X_{18}X_{201} - X_{19} - X_{174}X_{201} + X_{198} + X_{201} - X_{206}X_{207} \\[3pt]
r_{151} &= -4 - 4X_{0} - X_{0}^{2} - 2X_{0}X_{211} - 4X_{211} - X_{211}^{2} + X_{212} \\[3pt]
r_{152} &= -1 + X_{208}^{2} - X_{212}X_{230} \\[3pt]
r_{153} &= -X_{210}^{2} + X_{213} \\[3pt]
r_{154} &= -1 - X_{191}X_{213} + X_{209}^{2} + X_{213} \\[3pt]
r_{155} &= -X_{179}X_{214} - X_{208} + X_{209} \\[3pt]
r_{156} &= -2 - X_{0} - 4X_{0}X_{215} + X_{210} - 4X_{211}X_{215} - 8X_{215} \\[3pt]
r_{157} &= -2 - 2X_{22} - X_{22}^{2} + 2X_{22}X_{174} + 2X_{174} - X_{216} \\[3pt]
r_{158} &= -X_{0}X_{174} - 2X_{174} - X_{174}X_{211} + X_{218} \\[3pt]
r_{159} &= 2 + X_{0} + X_{0}X_{22} + 2X_{22} + X_{22}X_{211} - X_{23} + X_{208} + X_{211} - X_{216}X_{217} - X_{218} \\[3pt]
r_{160} &= -1 + 2X_{17}X_{174} - X_{18} - X_{219} \\[3pt]
r_{161} &= X_{0}X_{17} + 2X_{17} + X_{17}X_{211} - X_{24} + X_{208} - X_{218} - X_{219}X_{220} \\[3pt]
r_{162} &= 2 + X_{0} + X_{0}X_{18} + 2X_{18} + X_{18}X_{211} - X_{26} - X_{206}X_{221} + X_{208} + X_{211} - X_{218} \\[3pt]
r_{163} &= -X_{18}^{2} - 2X_{18}X_{225} - X_{225}^{2} + X_{226} \\[3pt]
r_{164} &= -1 + X_{222}^{2} - X_{226}X_{230} \\[3pt]
r_{165} &= -X_{224}^{2} + X_{227} \\[3pt]
r_{166} &= -1 - X_{191}X_{227} + X_{223}^{2} + X_{227} \\[3pt]
r_{167} &= -X_{179}X_{228} - X_{222} + X_{223} \\[3pt]
r_{168} &= -X_{18} - 4X_{18}X_{229} + X_{224} - 4X_{225}X_{229} \\[3pt]
r_{169} &= -X_{16}X_{18} - X_{16}X_{225} - X_{18}X_{201} - X_{201}X_{225} + X_{231} \\[3pt]
r_{170} &= -X_{18}X_{198} - X_{198}X_{225} + X_{232} \\[3pt]
r_{171} &= -X_{16}X_{222} - X_{201}X_{222} + X_{233} \\[3pt]
r_{172} &= -10 - 6X_{20} - X_{20}^{2} - 2X_{20}X_{21} - 2X_{20}X_{134} + 2X_{20}X_{174} - 6X_{21} - X_{21}^{2} - 2X_{21}X_{134} \\*
&\quad  + 2X_{21}X_{174} - 6X_{134} - X_{134}^{2} + 2X_{134}X_{174} + 6X_{174} - X_{234} \\[3pt]
r_{173} &= -1 - X_{20} - X_{20}X_{132} - X_{20}X_{133} + X_{20}X_{232} + X_{20}X_{233} - X_{21} - X_{21}X_{132} - X_{21}X_{133} + X_{21}X_{232} \\*
&\quad  + X_{21}X_{233} - 2X_{132} - X_{132}X_{135} - X_{133} - X_{133}X_{135} + X_{134}X_{232} + X_{134}X_{233} - X_{135} - X_{174}X_{232} \\*
&\quad  - X_{174}X_{233} + X_{198}X_{222} + X_{230}X_{231} + 3X_{232} + 3X_{233} - X_{234}X_{235} \\[3pt]
r_{174} &= -5 - 4X_{20} - X_{20}^{2} - 2X_{20}X_{21} - 2X_{20}X_{134} + 2X_{20}X_{174} - 4X_{21} - X_{21}^{2} - 2X_{21}X_{134} + 2X_{21}X_{174} \\*
&\quad  - 4X_{134} - X_{134}^{2} + 2X_{134}X_{174} + 4X_{174} - X_{236} \\[3pt]
r_{175} &= -1 + 2X_{16} + X_{16}X_{20} + X_{16}X_{21} + X_{16}X_{134} - X_{16}X_{174} + X_{20}X_{201} + X_{21}X_{201} - X_{132} - X_{133} \\*
&\quad  + X_{134}X_{201} - X_{174}X_{201} + X_{198} + 2X_{201} - X_{236}X_{237} \\[3pt]
r_{176} &= -X_{0}X_{18} - X_{0}X_{225} - 2X_{18} - X_{18}X_{211} - X_{211}X_{225} - 2X_{225} + X_{238} \\[3pt]
r_{177} &= -X_{18}X_{208} - X_{208}X_{225} + X_{239} \\[3pt]
r_{178} &= -X_{0}X_{222} - X_{211}X_{222} - 2X_{222} + X_{240} \\[3pt]
r_{179} &= -1 + X_{20}X_{239} + X_{20}X_{240} + X_{21}X_{239} + X_{21}X_{240} - X_{27} - X_{27}X_{137} - X_{27}X_{138} - X_{28} - X_{28}X_{137} \\*
&\quad  - X_{28}X_{138} + X_{134}X_{239} + X_{134}X_{240} - 2X_{137} - X_{137}X_{140} - X_{138} - X_{138}X_{140} - X_{140} - X_{174}X_{239} \\*
&\quad  - X_{174}X_{240} + X_{208}X_{222} + X_{230}X_{238} - X_{234}X_{241} + 3X_{239} + 3X_{240} \\[3pt]
r_{180} &= 3 + 2X_{0} + X_{0}X_{20} + X_{0}X_{21} + X_{0}X_{134} + 2X_{20} + X_{20}X_{211} + 2X_{21} + X_{21}X_{211} + 2X_{134} \\*
&\quad  + X_{134}X_{211} - X_{137} - X_{138} + X_{208} + 2X_{211} - X_{218} - X_{236}X_{242} \\[3pt]
r_{181} &= -16 - 8X_{0} - X_{0}^{2} - 2X_{0}X_{23} - 2X_{0}X_{246} - 8X_{23} - X_{23}^{2} - 2X_{23}X_{246} - 8X_{246} - X_{246}^{2} + X_{247} \\[3pt]
r_{182} &= -1 - X_{230}X_{247} + X_{243}^{2} \\[3pt]
r_{183} &= -X_{245}^{2} + X_{248} \\[3pt]
r_{184} &= -1 - X_{191}X_{248} + X_{244}^{2} + X_{248} \\[3pt]
r_{185} &= -X_{179}X_{249} - X_{243} + X_{244} \\[3pt]
r_{186} &= -4 - X_{0} - 4X_{0}X_{250} - X_{23} - 4X_{23}X_{250} + X_{245} - 4X_{246}X_{250} - 16X_{250} \\[3pt]
r_{187} &= 12 + 3X_{0} + X_{0}X_{20} + X_{0}X_{21} + X_{0}X_{134} - X_{0}X_{174} + 4X_{20} + X_{20}X_{23} + X_{20}X_{246} + 4X_{21} \\*
&\quad  + X_{21}X_{23} + X_{21}X_{246} + 3X_{23} + X_{23}X_{134} - X_{23}X_{174} - X_{29} - X_{29}X_{137} - X_{29}X_{138} + 4X_{134} \\*
&\quad  + X_{134}X_{246} - X_{137}X_{145} - X_{138}X_{145} - X_{142} - X_{145} - 4X_{174} - X_{174}X_{246} - X_{234}X_{251} + X_{243} \\*
&\quad  + 3X_{246} \\[3pt]
r_{188} &= -1 - 2X_{7} - X_{7}^{2} - 2X_{7}X_{255} - 2X_{255} - X_{255}^{2} + X_{256} \\[3pt]
r_{189} &= -1 - X_{230}X_{256} + X_{252}^{2} \\[3pt]
r_{190} &= -X_{254}^{2} + X_{257} \\[3pt]
r_{191} &= -1 - X_{191}X_{257} + X_{253}^{2} + X_{257} \\[3pt]
r_{192} &= -X_{179}X_{258} - X_{252} + X_{253} \\[3pt]
r_{193} &= -1 - X_{7} - 4X_{7}X_{259} + X_{254} - 4X_{255}X_{259} - 4X_{259} \\[3pt]
r_{194} &= 3 + 3X_{7} + X_{7}X_{20} + X_{7}X_{21} + X_{7}X_{134} - X_{7}X_{174} + X_{20} + X_{20}X_{255} + X_{21} + X_{21}X_{255} + X_{134} \\*
&\quad  + X_{134}X_{255} - X_{146}X_{149} - X_{147} - X_{174} - X_{174}X_{255} - X_{234}X_{260} + X_{252} + 3X_{255} \\[3pt]
r_{195} &= -1 + X_{0}X_{230} + X_{20}X_{208} + X_{20}X_{218} + X_{21}X_{208} + X_{21}X_{218} + X_{134}X_{208} + X_{134}X_{218} - X_{147} \\*
&\quad  - X_{148} - X_{174}X_{218} + 2X_{208} + X_{211}X_{230} + 2X_{218} + 2X_{230} - X_{236}X_{261} \\[3pt]
r_{196} &= -1 - 2X_{8} - X_{8}^{2} - 2X_{8}X_{265} - 2X_{265} - X_{265}^{2} + X_{266} \\[3pt]
r_{197} &= -1 - X_{230}X_{266} + X_{262}^{2} \\[3pt]
r_{198} &= -X_{264}^{2} + X_{267} \\[3pt]
r_{199} &= -1 - X_{191}X_{267} + X_{263}^{2} + X_{267} \\[3pt]
r_{200} &= -X_{179}X_{268} - X_{262} + X_{263} \\[3pt]
r_{201} &= -1 - X_{8} - 4X_{8}X_{269} + X_{264} - 4X_{265}X_{269} - 4X_{269} \\[3pt]
r_{202} &= 3 + 3X_{8} + X_{8}X_{20} + X_{8}X_{21} + X_{8}X_{134} - X_{8}X_{174} + X_{20} + X_{20}X_{265} + X_{21} + X_{21}X_{265} + X_{134} \\*
&\quad  + X_{134}X_{265} - X_{149}X_{150} - X_{151} - X_{174} - X_{174}X_{265} - X_{234}X_{270} + X_{262} + 3X_{265} \\[3pt]
r_{203} &= -1 - 2X_{9} - X_{9}^{2} - 2X_{9}X_{274} - 2X_{274} - X_{274}^{2} + X_{275} \\[3pt]
r_{204} &= -1 - X_{230}X_{275} + X_{271}^{2} \\[3pt]
r_{205} &= -X_{273}^{2} + X_{276} \\[3pt]
r_{206} &= -1 - X_{191}X_{276} + X_{272}^{2} + X_{276} \\[3pt]
r_{207} &= -X_{179}X_{277} - X_{271} + X_{272} \\[3pt]
r_{208} &= -1 - X_{9} - 4X_{9}X_{278} + X_{273} - 4X_{274}X_{278} - 4X_{278} \\[3pt]
r_{209} &= 3 + 3X_{9} + X_{9}X_{20} + X_{9}X_{21} + X_{9}X_{134} - X_{9}X_{174} + X_{20} + X_{20}X_{274} + X_{21} + X_{21}X_{274} + X_{134} \\*
&\quad  + X_{134}X_{274} - X_{149}X_{153} - X_{154} - X_{174} - X_{174}X_{274} - X_{234}X_{279} + X_{271} + 3X_{274} \\[3pt]
r_{210} &= -1 - 2X_{10} - X_{10}^{2} - 2X_{10}X_{283} - 2X_{283} - X_{283}^{2} + X_{284} \\[3pt]
r_{211} &= -1 - X_{230}X_{284} + X_{280}^{2} \\[3pt]
r_{212} &= -X_{282}^{2} + X_{285} \\[3pt]
r_{213} &= -1 - X_{191}X_{285} + X_{281}^{2} + X_{285} \\[3pt]
r_{214} &= -X_{179}X_{286} - X_{280} + X_{281} \\[3pt]
r_{215} &= -1 - X_{10} - 4X_{10}X_{287} + X_{282} - 4X_{283}X_{287} - 4X_{287} \\[3pt]
r_{216} &= 3 + 3X_{10} + X_{10}X_{20} + X_{10}X_{21} + X_{10}X_{134} - X_{10}X_{174} + X_{20} + X_{20}X_{283} + X_{21} + X_{21}X_{283} \\*
&\quad  + X_{134} + X_{134}X_{283} - X_{149}X_{156} - X_{157} - X_{174} - X_{174}X_{283} - X_{234}X_{288} + X_{280} + 3X_{283} \\[3pt]
r_{217} &= -1 - 2X_{11} - X_{11}^{2} - 2X_{11}X_{292} - 2X_{292} - X_{292}^{2} + X_{293} \\[3pt]
r_{218} &= -1 - X_{230}X_{293} + X_{289}^{2} \\[3pt]
r_{219} &= -X_{291}^{2} + X_{294} \\[3pt]
r_{220} &= -1 - X_{191}X_{294} + X_{290}^{2} + X_{294} \\[3pt]
r_{221} &= -X_{179}X_{295} - X_{289} + X_{290} \\[3pt]
r_{222} &= -1 - X_{11} - 4X_{11}X_{296} + X_{291} - 4X_{292}X_{296} - 4X_{296} \\[3pt]
r_{223} &= 3 + 3X_{11} + X_{11}X_{20} + X_{11}X_{21} + X_{11}X_{134} - X_{11}X_{174} + X_{20} + X_{20}X_{292} + X_{21} + X_{21}X_{292} \\*
&\quad  + X_{134} + X_{134}X_{292} - X_{159}X_{162} - X_{160} - X_{174} - X_{174}X_{292} - X_{234}X_{297} + X_{289} + 3X_{292} \\[3pt]
r_{224} &= -9 - 6X_{2} - X_{2}^{2} - 2X_{2}X_{4} - 2X_{2}X_{301} - 6X_{4} - X_{4}^{2} - 2X_{4}X_{301} - 6X_{301} - X_{301}^{2} + X_{302} \\[3pt]
r_{225} &= -1 - X_{230}X_{302} + X_{298}^{2} \\[3pt]
r_{226} &= -X_{300}^{2} + X_{303} \\[3pt]
r_{227} &= -1 - X_{191}X_{303} + X_{299}^{2} + X_{303} \\[3pt]
r_{228} &= -X_{179}X_{304} - X_{298} + X_{299} \\[3pt]
r_{229} &= -3 - X_{2} - 4X_{2}X_{305} - X_{4} - 4X_{4}X_{305} + X_{300} - 4X_{301}X_{305} - 12X_{305} \\[3pt]
r_{230} &= 5 + 2X_{2} + X_{2}X_{20} + X_{2}X_{21} + X_{2}X_{134} - X_{2}X_{174} + 2X_{4} + X_{4}X_{20} + X_{4}X_{21} + X_{4}X_{134} \\*
&\quad  - X_{4}X_{174} + 3X_{20} + X_{20}X_{301} + 3X_{21} + X_{21}X_{301} + 3X_{134} + X_{134}X_{301} - X_{160} - X_{161} - 3X_{174} \\*
&\quad  - X_{174}X_{301} - X_{236}X_{306} + X_{298} + 2X_{301} \\[3pt]
r_{231} &= -1 - 2X_{12} - X_{12}^{2} - 2X_{12}X_{310} - 2X_{310} - X_{310}^{2} + X_{311} \\[3pt]
r_{232} &= -1 - X_{230}X_{311} + X_{307}^{2} \\[3pt]
r_{233} &= -X_{309}^{2} + X_{312} \\[3pt]
r_{234} &= -1 - X_{191}X_{312} + X_{308}^{2} + X_{312} \\[3pt]
r_{235} &= -X_{179}X_{313} - X_{307} + X_{308} \\[3pt]
r_{236} &= -1 - X_{12} - 4X_{12}X_{314} + X_{309} - 4X_{310}X_{314} - 4X_{314} \\[3pt]
r_{237} &= 3 + 3X_{12} + X_{12}X_{20} + X_{12}X_{21} + X_{12}X_{134} - X_{12}X_{174} + X_{20} + X_{20}X_{310} + X_{21} + X_{21}X_{310} \\*
&\quad  + X_{134} + X_{134}X_{310} - X_{162}X_{163} - X_{164} - X_{174} - X_{174}X_{310} - X_{234}X_{315} + X_{307} + 3X_{310} \\[3pt]
r_{238} &= -1 - 2X_{13} - X_{13}^{2} - 2X_{13}X_{319} - 2X_{319} - X_{319}^{2} + X_{320} \\[3pt]
r_{239} &= -1 - X_{230}X_{320} + X_{316}^{2} \\[3pt]
r_{240} &= -X_{318}^{2} + X_{321} \\[3pt]
r_{241} &= -1 - X_{191}X_{321} + X_{317}^{2} + X_{321} \\[3pt]
r_{242} &= -X_{179}X_{322} - X_{316} + X_{317} \\[3pt]
r_{243} &= -1 - X_{13} - 4X_{13}X_{323} + X_{318} - 4X_{319}X_{323} - 4X_{323} \\[3pt]
r_{244} &= 3 + 3X_{13} + X_{13}X_{20} + X_{13}X_{21} + X_{13}X_{134} - X_{13}X_{174} + X_{20} + X_{20}X_{319} + X_{21} + X_{21}X_{319} \\*
&\quad  + X_{134} + X_{134}X_{319} - X_{166}X_{169} - X_{167} - X_{174} - X_{174}X_{319} - X_{234}X_{324} + X_{316} + 3X_{319} \\[3pt]
r_{245} &= -1 - 2X_{1} - X_{1}^{2} - 2X_{1}X_{328} - 2X_{328} - X_{328}^{2} + X_{329} \\[3pt]
r_{246} &= -1 - X_{230}X_{329} + X_{325}^{2} \\[3pt]
r_{247} &= -X_{327}^{2} + X_{330} \\[3pt]
r_{248} &= -1 - X_{191}X_{330} + X_{326}^{2} + X_{330} \\[3pt]
r_{249} &= -X_{179}X_{331} - X_{325} + X_{326} \\[3pt]
r_{250} &= -1 - X_{1} - 4X_{1}X_{332} + X_{327} - 4X_{328}X_{332} - 4X_{332} \\[3pt]
r_{251} &= 1 + 2X_{1} + X_{1}X_{20} + X_{1}X_{21} + X_{1}X_{134} - X_{1}X_{174} + X_{20} + X_{20}X_{328} + X_{21} + X_{21}X_{328} + X_{134} \\*
&\quad  + X_{134}X_{328} - X_{167} - X_{168} - X_{174} - X_{174}X_{328} - X_{236}X_{333} + X_{325} + 2X_{328} \\[3pt]
r_{252} &= -1 - 2X_{14} - X_{14}^{2} - 2X_{14}X_{337} - 2X_{337} - X_{337}^{2} + X_{338} \\[3pt]
r_{253} &= -1 - X_{230}X_{338} + X_{334}^{2} \\[3pt]
r_{254} &= -X_{336}^{2} + X_{339} \\[3pt]
r_{255} &= -1 - X_{191}X_{339} + X_{335}^{2} + X_{339} \\[3pt]
r_{256} &= -X_{179}X_{340} - X_{334} + X_{335} \\[3pt]
r_{257} &= -1 - X_{14} - 4X_{14}X_{341} + X_{336} - 4X_{337}X_{341} - 4X_{341} \\[3pt]
r_{258} &= 3 + 3X_{14} + X_{14}X_{20} + X_{14}X_{21} + X_{14}X_{134} - X_{14}X_{174} + X_{20} + X_{20}X_{337} + X_{21} + X_{21}X_{337} \\*
&\quad  + X_{134} + X_{134}X_{337} - X_{169}X_{170} - X_{171} - X_{174} - X_{174}X_{337} - X_{234}X_{342} + X_{334} + 3X_{337} \\[3pt]
r_{259} &= -X_{189}X_{202} + X_{343} \\[3pt]
r_{260} &= -X_{212}X_{343} + X_{344} \\[3pt]
r_{261} &= -X_{226}X_{344} + X_{345} \\[3pt]
r_{262} &= -X_{247}X_{345} + X_{346} \\[3pt]
r_{263} &= -X_{256}X_{346} + X_{347} \\[3pt]
r_{264} &= -X_{266}X_{347} + X_{348} \\[3pt]
r_{265} &= -X_{275}X_{348} + X_{349} \\[3pt]
r_{266} &= -X_{284}X_{349} + X_{350} \\[3pt]
r_{267} &= -X_{293}X_{350} + X_{351} \\[3pt]
r_{268} &= -X_{302}X_{351} + X_{352} \\[3pt]
r_{269} &= -X_{311}X_{352} + X_{353} \\[3pt]
r_{270} &= -X_{320}X_{353} + X_{354} \\[3pt]
r_{271} &= -X_{329}X_{354} + X_{355} \\[3pt]
r_{272} &= -X_{338}X_{355} + X_{356} \\[3pt]
r_{273} &= X_{180} - X_{356} - X_{356}X_{357} \\[3pt]
r_{274} &= -8388654 - 4194309X_{0} - 2X_{1} - 4X_{2} - 2X_{4} - X_{7} - X_{8} - X_{9} - X_{10} - X_{11} - X_{12} - X_{13} - X_{14} \\*
&\quad  - 2X_{16} - X_{17} - 3X_{18} - X_{19} - 3X_{20} - X_{20}X_{132} - X_{20}X_{133} - 3X_{21} - X_{21}X_{132} - X_{21}X_{133} - X_{22} \\*
&\quad  - 2X_{23} - X_{24} - X_{26} - X_{27} - X_{27}X_{137} - X_{27}X_{138} - X_{28} - X_{28}X_{137} - X_{28}X_{138} - X_{29} - X_{29}X_{137} \\*
&\quad  - X_{29}X_{138} - 3X_{132} - X_{132}X_{135} - 2X_{133} - X_{133}X_{135} - 2X_{134} - X_{135} - 3X_{137} - X_{137}X_{140} \\*
&\quad  - X_{137}X_{145} - 2X_{138} - X_{138}X_{140} - X_{138}X_{145} - X_{140} - X_{142} - X_{145} - X_{146}X_{149} - 2X_{147} - X_{148} \\*
&\quad  - X_{149}X_{150} - X_{149}X_{153} - X_{149}X_{156} - X_{151} - X_{154} - X_{157} - X_{159}X_{162} - 2X_{160} - X_{161} \\*
&\quad  - X_{162}X_{163} - X_{164} - X_{166}X_{169} - 2X_{167} - X_{168} - X_{169}X_{170} - X_{171} + X_{173}
\end{align*}
\endgroup


\end{document}
